\chapter{Ciała rozkładu i grupy Galois}

\section{Ciała rozkładu, normalność i rozdzielczość}

\begin{definicja}
Ciałem rozkładu niezerowego wielomianu $f\in K[T]$ jest rozszerzenie $L=K(\alpha_1,\dots,\alpha_r)$ generowane przez wszystkie jego różne zera, w którym $f$ rozkłada się na czynniki liniowe. Wielomian jest \emph{rozdzielczy}, jeśli nie ma zer wielokrotnych.
\end{definicja}

\begin{twierdzenie}
Ciało rozkładu istnieje, jest skończone nad $K$ i jest jednoznaczne z dokładnością do izomorfizmu nad $K$.
\end{twierdzenie}
\begin{proof}
Jeżeli wielomian nie jest już iloczynem czynników liniowych, dołączamy zero jednego z jego nierozkładalnych czynników, korzystając z twierdzenia 1.3. W nowym ciele wydzielamy czynnik liniowy. Indukcja po stopniu pozostałej części kończy się po skończenie wielu krokach. Każdy krok jest skończony, więc twierdzenie o wieży daje skończoność całego rozszerzenia.

Niech $L$ i $M$ będą dwoma ciałami rozkładu. Zanurzamy $L$ w domknięciu algebraicznym $M$, przedłużając identyczność $K$. Obrazy wszystkich różnych zer $f$ są różnymi zerami $f$, zatem są dokładnie wszystkimi zerami w $M$. Obraz $L$ jest ciałem przez nie generowanym, czyli $M$. Otrzymane zanurzenie jest więc izomorfizmem $L\simeq M$ nad $K$.
\end{proof}

Rozszerzenie algebraiczne jest \emph{normalne}, jeśli każdy nierozkładalny wielomian nad bazą mający w nim zero rozkłada się tam całkowicie. Jest \emph{rozdzielcze}, jeśli wszystkie wielomiany minimalne jego elementów są rozdzielcze. W charakterystyce zero druga własność zachodzi automatycznie, jak pokazano w dowodzie liczby zanurzeń.

\begin{lemat}[Trzy opisy normalności]
Dla skończonego rozszerzenia $L/K$ zawartego w $\overline K$ równoważne są:
\begin{enumerate}
\item $L/K$ jest normalne;
\item każde $K$-zanurzenie $\sigma:L\to\overline K$ spełnia $\sigma(L)=L$;
\item $L$ jest ciałem rozkładu pewnego wielomianu nad $K$.
\end{enumerate}
\end{lemat}
\begin{proof}
Jeżeli $L/K$ jest normalne, obraz każdego $\alpha\in L$ jest zerem jego wielomianu minimalnego, więc należy do $L$. Zanurzenie $L\to L$ jest surjektywne, bo jest różnowartościowym operatorem $K$-liniowym na przestrzeni skończonego wymiaru. To daje (1)$\Rightarrow$(2).

Załóżmy (2). Jeśli $\alpha\in L$ i $\beta$ jest dowolnym zerem $m_{\alpha,K}$, to zanurzenie $K(\alpha)\to\overline K$, $\alpha\mapsto\beta$, przedłuża się do $L$. Warunek (2) daje $\beta\in L$. Zatem $L/K$ jest normalne.

Przy (1) wybieramy generatory $\alpha_1,\dots,\alpha_r$ ciała $L/K$. Iloczyn ich wielomianów minimalnych rozkłada się w $L$, a jego zera generują dokładnie $L$. Otrzymujemy (3). Przy (3) każde zanurzenie permutuje cały skończony zbiór zer wielomianu definiującego ciało rozkładu; stąd jego obrazem jest $L$, co daje (2).
\end{proof}

\begin{definicja}
Skończone normalne i rozdzielcze rozszerzenie nazywamy \emph{rozszerzeniem Galois}. Jego grupa Galois to
\[
\Gal(L/K)=\Aut_K(L)=\{\sigma:L\xrightarrow{\sim}L:\sigma|_K=\id_K\}.
\]
Grupa Galois wielomianu $f$ oznacza grupę jego ciała rozkładu nad ciałem współczynników.
\end{definicja}

\begin{twierdzenie}[Liczba automorfizmów]
Dla skończonego $L/K$ w charakterystyce zero zachodzi $|\Aut_K(L)|\le[L:K]$. Równość zachodzi wtedy i tylko wtedy, gdy $L/K$ jest Galois.
\end{twierdzenie}
\begin{proof}
Automorfizmy są szczególnymi zanurzeniami, których jest $[L:K]$. Równość oznacza, że każde zanurzenie ma obraz $L$. Jest to dokładnie normalność z lematu 2.3; rozdzielczość jest automatyczna.
\end{proof}

Jeśli $f$ ma $n$ różnych zer, grupa jego ciała rozkładu działa na nich wiernie: automorfizm ustalający wszystkie generatory jest identycznością. Możemy więc traktować ją jako podgrupę $S_n$. Gdy $f$ jest nierozkładalny, działanie jest przechodnie. Istotnie, zanurzenie wysyłające jedno zero na dowolne inne przedłuża się do ciała rozkładu, a normalność zamienia to przedłużenie w automorfizm.

\section{Trzy obliczenia grup Galois}

\begin{przyklad}[Rozszerzenie kwadratowe]
Jeśli $d\in\Q$ nie jest kwadratem, to $\Q(\sqrt d)/\Q$ ma stopień $2$ i jest ciałem rozkładu $T^2-d$. Jego automorfizmy to identyczność oraz $a+b\sqrt d\mapsto a-b\sqrt d$. Grupa jest izomorficzna z $C_2$.
\end{przyklad}

\begin{przyklad}[Ciało rozkładu wielomianu sześciennego]
Niech $\alpha=\sqrt[3]{2}$ i $\omega=e^{2\pi i/3}$. Wielomian $T^3-2$ nie ma zera wymiernego: dla nieskracalnego $a/b$ równanie $a^3=2b^3$ przeczy porównaniu wykładników liczby pierwszej $2$. Wielomian stopnia $3$ bez zera wymiernego jest nierozkładalny, ponieważ w każdym jego nietrywialnym rozkładzie jeden czynnik musiałby być liniowy.

Ciało rozkładu to $L=\Q(\alpha,\omega)$. Mamy $[\Q(\alpha):\Q]=3$ oraz $[L:\Q(\alpha)]=2$: element $\omega$ spełnia $T^2+T+1=0$ i nie należy do rzeczywistego ciała $\Q(\alpha)$. Zatem $[L:\Q]=6$. Grupa Galois ma sześć elementów i zanurza się w $S_3$, więc jest całą $S_3$. W szczególności istnieją automorfizmy
\[
\sigma(\alpha)=\omega\alpha,\quad\sigma(\omega)=\omega,
\qquad\tau(\alpha)=\alpha,\quad\tau(\omega)=\omega^2.
\]
Spełniają $\sigma^3=\tau^2=1$ i $\tau\sigma\tau=\sigma^{-1}$. Istnienie $\sigma$ wynika tu z już obliczonej pełnej grupy permutacji, a nie tylko z formalnego przypisania obrazów generatorom.
\end{przyklad}

\begin{przyklad}[Rozszerzenie dwukwadratowe]
Dla $L=\Q(\sqrt2,\sqrt3)$ sprawdzamy najpierw stopień. Gdyby $\sqrt3=a+b\sqrt2$ z $a,b\in\Q$, to po podniesieniu do kwadratu $3=a^2+2b^2+2ab\sqrt2$. Niewymierność $\sqrt2$ daje $ab=0$. Otrzymalibyśmy $a^2=3$ albo $b^2=3/2$, a żadna z tych liczb nie jest kwadratem wymiernym (wykładniki liczb pierwszych w kwadracie są parzyste). Stąd $[L:\Q]=4$. Jest to ciało rozkładu $(T^2-2)(T^2-3)$.

Każdy automorfizm wybiera znaki obrazów obu pierwiastków. Są co najwyżej cztery takie wybory i dokładnie cztery automorfizmy, więc wszystkie wybory występują. Zmiany znaków komutują i mają rząd $2$, toteż $\Gal(L/\Q)\simeq C_2\times C_2$.
\end{przyklad}


